\section{An infinite incidence ladder}

Define two admissible cyclic families
\begin{equation}\label{eq:families}
 A_m=0^{m-2}21\quad(m\geq3),
 \qquad
 B_m=0^{m-3}231\quad(m\geq4).
\end{equation}
Every displayed transition is allowed by \(A\), including the cyclic return
\(1\to0\).  Each word contains exactly one symbol \(2\), so it cannot be a
proper power.  Thus all members are primitive.

Write \(e_w\) for the coordinate vector of a block \(w\).

\begin{lemma}[common insertion row]\label{lem:insertion}
For every \(m\geq3\),
\begin{align}
 \Nblk_m(A_{m+1})-\Nblk_m(A_m)
 &=e_{0^{m-1}2}+e_{10^{m-1}}-e_{10^{m-2}2},\label{eq:A-insertion}\\
 \Nblk_m(B_{m+2})-\Nblk_m(B_{m+1})
 &=e_{0^{m-1}2}+e_{10^{m-1}}-e_{10^{m-2}2}.\label{eq:B-insertion}
\end{align}
\end{lemma}

\begin{proof}
In each comparison one extra zero is inserted at the beginning of the
unique maximal zero run.  List the cyclic width-\(m\) windows by their start
relative to the unique terminal symbol \(1\).  Windows that see the same
part of the zero run cancel bijectively.  The longer word adds the windows
\(0^{m-1}2\) and \(10^{m-1}\), while the shorter word has the unmatched
window \(10^{m-2}2\).  The symbols after \(2\) distinguish the two families
but lie in the portion common to both sides of each comparison, so the same
three survivors remain.  This gives both identities.
\end{proof}

\begin{theorem}[all-width block-incidence ladder]\label{thm:ladder}
For every integer \(m\geq3\),
\begin{equation}\label{eq:ladder}
 \Nblk_m(A_m)+\Nblk_m(B_{m+2})
 =\Nblk_m(A_{m+1})+\Nblk_m(B_{m+1}).
\end{equation}
\end{theorem}

\begin{proof}
Subtract \eqref{eq:A-insertion} from \eqref{eq:B-insertion} and rearrange.
\end{proof}

The cases \(m=3\) and \(m=4\) recover respectively the P55 relation and its
first extension:
\begin{align*}
 \Nblk_3(A_3)+\Nblk_3(B_5)&=\Nblk_3(A_4)+\Nblk_3(B_4),\\
 \Nblk_4(A_4)+\Nblk_4(B_6)&=\Nblk_4(A_5)+\Nblk_4(B_5).
\end{align*}
The exact program checks the formula independently through \(m=64\), but
that finite check is evidence for the implementation, not the proof.
